\PassOptionsToPackage{unicode}{hyperref}
\PassOptionsToPackage{hyphens}{url}
\PassOptionsToPackage{dvipsnames,svgnames,x11names}{xcolor}
\documentclass[
  spanish,
  11pt,
  a4paper,
  DIV=11,
  numbers=noendperiod]{scrartcl}

\usepackage{xcolor}
\usepackage[top=24mm,bottom=24mm,left=24mm,right=24mm,headheight=15pt,headsep=8mm]{geometry}
\usepackage{amsmath,amssymb}
\usepackage{iftex}
\ifPDFTeX
  \usepackage[T1]{fontenc}
  \usepackage[utf8]{inputenc}
  \usepackage{textcomp}
\else
  \usepackage{unicode-math}
  \defaultfontfeatures{Scale=MatchLowercase}
  \defaultfontfeatures[\rmfamily]{Ligatures=TeX,Scale=1}
\fi
\usepackage{lmodern}
\IfFileExists{upquote.sty}{\usepackage{upquote}}{}
\IfFileExists{microtype.sty}{
  \usepackage[]{microtype}
  \UseMicrotypeSet[protrusion]{basicmath}
}{}
\usepackage{graphicx}
\usepackage{float}
\usepackage[bidi=basic]{babel}
\let\LanguageShortHands\languageshorthands
\def\languageshorthands#1{}
\setlength{\emergencystretch}{3em}
\setcounter{secnumdepth}{5}
\providecommand{\tightlist}{\setlength{\itemsep}{0pt}\setlength{\parskip}{0pt}}

\newcommand{\practicaPdfCoverImage}{../resources/practica3/images/teaser.jpg}
% PDF style for Practica 9.
% A modest university handout look: classical type, clear hierarchy,
% restrained color, and durable code/callout styling.

\usepackage{graphicx}
\usepackage[export]{adjustbox}
\usepackage{fontspec}
\usepackage{microtype}
\usepackage{scrlayer-scrpage}
\usepackage{tikz}
\usepackage{fvextra}
\usepackage{caption}
\usepackage{subcaption}
\usepackage[skins,breakable]{tcolorbox}

\providecommand{\practicaPdfCoverImage}{}

\definecolor{uzaInk}{HTML}{172033}
\definecolor{uzaMuted}{HTML}{5F6B7A}
\definecolor{uzaBlue}{HTML}{1F5F8B}
\definecolor{uzaBlueDark}{HTML}{153E5C}
\definecolor{uzaCyan}{HTML}{EAF4F7}
\definecolor{uzaLine}{HTML}{D7E2EA}
\definecolor{uzaCodeBg}{HTML}{F5F7F9}
\definecolor{uzaTip}{HTML}{237A57}
\definecolor{uzaWarn}{HTML}{A65F00}

\IfFontExistsTF{TeX Gyre Pagella}{
  \setmainfont{TeX Gyre Pagella}
}{}
\IfFontExistsTF{TeX Gyre Heros}{
  \setsansfont{TeX Gyre Heros}
}{}
\IfFontExistsTF{TeX Gyre Cursor}{
  \setmonofont{TeX Gyre Cursor}[Scale=0.86]
}{
  \setmonofont{Latin Modern Mono}[Scale=0.88]
}

\KOMAoptions{parskip=half}
\setlength{\parindent}{0pt}
\setlength{\parskip}{0.55em plus 0.15em minus 0.08em}
\linespread{1.035}

\addtokomafont{disposition}{\sffamily\color{uzaBlueDark}}
\addtokomafont{section}{\LARGE}
\addtokomafont{subsection}{\Large}
\addtokomafont{subsubsection}{\normalsize\bfseries}
\RedeclareSectionCommand[beforeskip=2.2em plus 0.4em minus 0.2em,afterskip=1.0em]{section}
\RedeclareSectionCommand[beforeskip=1.2em,afterskip=0.55em]{subsection}
\RedeclareSectionCommand[beforeskip=1.0em,afterskip=0.35em]{subsubsection}

\clearpairofpagestyles
\automark[section]{section}
\ihead{\sffamily\footnotesize\textcolor{uzaMuted}{Informática (30717)}}
\ohead{\sffamily\footnotesize\textcolor{uzaMuted}{\headmark}}
\cfoot{\sffamily\small\textcolor{uzaMuted}{\pagemark}}
\setkomafont{pageheadfoot}{\normalfont}
\KOMAoptions{headsepline=0.35pt}
\addtokomafont{headsepline}{\color{uzaLine}}

\makeatletter
\AtBeginDocument{%
  \hypersetup{
    linkcolor=uzaBlueDark,
    urlcolor=uzaBlue,
    citecolor=uzaBlueDark
  }%
  \@ifundefined{Highlighting}{%
    \DefineVerbatimEnvironment{Highlighting}{Verbatim}{
      commandchars=\\\{\},
      fontsize=\small,
      breaklines=true,
      breakanywhere=true
    }%
  }{%
    \RecustomVerbatimEnvironment{Highlighting}{Verbatim}{
      commandchars=\\\{\},
      fontsize=\small,
      breaklines=true,
      breakanywhere=true
    }%
  }%
  \@ifundefined{Shaded}{%
    \newenvironment{Shaded}{%
      \begin{tcolorbox}[
        enhanced,
        breakable,
        sharp corners,
        boxrule=0.25pt,
        colback=uzaCodeBg,
        colframe=uzaLine,
        borderline west={1.4pt}{0pt}{uzaBlue},
        left=2mm,
        right=2mm,
        top=1.4mm,
        bottom=1.4mm,
        before skip=0.75em,
        after skip=0.75em
      ]%
    }{%
      \end{tcolorbox}%
    }%
  }{%
    \renewenvironment{Shaded}{%
      \begin{tcolorbox}[
        enhanced,
        breakable,
        sharp corners,
        boxrule=0.25pt,
        colback=uzaCodeBg,
        colframe=uzaLine,
        borderline west={1.4pt}{0pt}{uzaBlue},
        left=2mm,
        right=2mm,
        top=1.4mm,
        bottom=1.4mm,
        before skip=0.75em,
        after skip=0.75em
      ]%
    }{%
      \end{tcolorbox}%
    }%
  }%
}
\makeatother

\newcommand{\PracticeCode}[1]{%
  \begin{Shaded}%
  \VerbatimInput[
    fontsize=\small,
    breaklines=true,
    breakanywhere=true
  ]{#1}%
  \end{Shaded}%
}

\makeatletter
\renewcommand{\maketitle}{%
  \begin{titlepage}
    \thispagestyle{empty}
    \begin{tikzpicture}[remember picture,overlay]
      \fill[uzaCyan] (current page.north west) rectangle ([yshift=-72mm]current page.north east);
      \fill[uzaBlue] (current page.north west) rectangle ([xshift=7mm]current page.south west);
    \end{tikzpicture}
    \vspace*{12mm}
    {\sffamily\small\bfseries\textcolor{uzaBlue}{INFORMÁTICA (30717)}}\par
    \vspace{5mm}
    {\sffamily\bfseries\fontsize{30}{35}\selectfont\textcolor{uzaInk}{\@title}}\par
    \vspace{3mm}
    {\sffamily\Large\textcolor{uzaMuted}{Grado en Estudios en Arquitectura}}\par
    \vspace{8mm}
    {\color{uzaBlue}\rule{\textwidth}{0.6pt}}\par
    \if\relax\detokenize\expandafter{\practicaPdfCoverImage}\relax
    \else
      \vspace{5mm}
      {\centering\includegraphics[width=\textwidth,height=78mm,keepaspectratio]{\practicaPdfCoverImage}\par}
    \fi
    \vfill
    {\sffamily\large\textcolor{uzaBlueDark}{Universidad de Zaragoza}}\par
    \vspace{2mm}
    {\sffamily\small\textcolor{uzaMuted}{Material de prácticas del curso 2025--2026}}\par
  \end{titlepage}
}
\makeatother

\captionsetup{
  font={small,sf},
  labelfont={bf,color=uzaBlueDark},
  textfont={color=uzaInk},
  justification=centering,
  singlelinecheck=false,
  skip=6pt
}
\captionsetup[subfigure]{
  font={small,sf},
  labelfont={bf,color=uzaBlueDark},
  justification=centering,
  singlelinecheck=true,
  skip=4pt
}

\newcommand{\PracticeCredit}[1]{\textcolor{uzaMuted}{#1}}
\newcommand{\PracticeImagePair}[6]{%
  \begin{figure}[tbp]
    \centering
    \begin{minipage}[t]{0.485\linewidth}
      \includegraphics[width=\linewidth]{#3}
      \vspace{2mm}
      \centering{\sffamily\footnotesize #4}
    \end{minipage}\hfill
    \begin{minipage}[t]{0.485\linewidth}
      \includegraphics[width=\linewidth]{#5}
      \vspace{2mm}
      \centering{\sffamily\footnotesize #6}
    \end{minipage}
    \caption{\label{#1}#2}
  \end{figure}
}

\renewcommand{\arraystretch}{1.16}
\setlength{\tabcolsep}{6pt}
\setkeys{Gin}{center}

% Quarto callout colors are referenced by generated tcolorbox options.
% Apply them at document start so they win over Quarto's defaults.
\AtBeginDocument{%
  \definecolor{quarto-callout-note-color}{HTML}{1F5F8B}%
  \definecolor{quarto-callout-note-color-frame}{HTML}{1F5F8B}%
  \definecolor{quarto-callout-tip-color}{HTML}{237A57}%
  \definecolor{quarto-callout-tip-color-frame}{HTML}{237A57}%
  \definecolor{quarto-callout-warning-color}{HTML}{A65F00}%
  \definecolor{quarto-callout-warning-color-frame}{HTML}{A65F00}%
  \definecolor{quarto-callout-important-color}{HTML}{9B2C2C}%
  \definecolor{quarto-callout-important-color-frame}{HTML}{9B2C2C}%
  \definecolor{quarto-callout-caution-color}{HTML}{B45309}%
  \definecolor{quarto-callout-caution-color-frame}{HTML}{B45309}%
}

\KOMAoption{captions}{tableheading}

\usepackage{bookmark}
\IfFileExists{xurl.sty}{\usepackage{xurl}}{}
\urlstyle{same}
\hypersetup{
  pdftitle={Práctica 3: Resolución de problemas en Java (I)},
  pdflang={es},
  colorlinks=true,
  linkcolor={uzaBlueDark},
  filecolor={Maroon},
  citecolor={Blue},
  urlcolor={uzaBlue},
  pdfcreator={LaTeX}
}

\title{Práctica 3: Resolución de problemas en Java (I)}
\author{}
\date{}

\begin{document}
\maketitle

\renewcommand*\contentsname{Tabla de contenidos}
{
\hypersetup{linkcolor=}
\setcounter{tocdepth}{2}
\tableofcontents
}

\clearpage

\section{Objetivos de la práctica}\label{objetivos-de-la-practica}

Los objetivos de esta práctica son los siguientes:

\begin{itemize}
\tightlist
\item Desarrollar algoritmos en Java.
\item Ejecutar algoritmos en el entorno Eclipse.
\item Utilizar expresiones condicionales para distinguir distintos casos de un problema.
\end{itemize}

\section{Raíces de una ecuación de segundo grado}\label{raices-de-una-ecuacion-de-segundo-grado}

Una ecuación de segundo grado tiene la forma:

\[
ax^2 + bx + c = 0
\]

donde \(a\), \(b\) y \(c\) son coeficientes reales, y \(x\) es la variable que queremos resolver.

A diferencia de las ecuaciones de primer grado, las ecuaciones de segundo grado pueden tener \textbf{0, 1 o 2 soluciones reales} dependiendo de los valores de los coeficientes.

Por ejemplo, los siguientes coeficientes generan diferentes números de soluciones:

\begin{itemize}
\tightlist
\item \(a=1\), \(b=-3\), \(c=2\): dos soluciones reales (\(x=1\) y \(x=2\)).
\item \(a=1\), \(b=-2\), \(c=1\): una solución real (\(x=1\)).
\item \(a=1\), \(b=0\), \(c=1\): no hay soluciones reales.
\end{itemize}

\begin{figure}[H]
\centering
\begin{tikzpicture}[x=1cm,y=1cm, every node/.style={font=\sffamily\small}]
  \foreach \shift/\title/\formula/\roots in {
    0/{$\Delta>0$}/{Dos cortes}/{dos raíces reales},
    5/{$\Delta=0$}/{Tangencia}/{una raíz real},
    10/{$\Delta<0$}/{Sin cortes}/{sin raíces reales}
  } {
    \begin{scope}[xshift=\shift cm]
      \draw[rounded corners=2pt, color=uzaLine, fill=uzaCodeBg] (-2.05,-1.35) rectangle (2.05,2.05);
      \draw[->, color=uzaMuted] (-1.65,0) -- (1.72,0) node[right] {$x$};
      \draw[->, color=uzaMuted] (0,-0.95) -- (0,1.55) node[above] {$y$};
      \node[anchor=west, color=uzaBlueDark, font=\sffamily\bfseries\small] at (-1.8,1.7) {\title};
      \node[anchor=west, color=uzaMuted, font=\sffamily\scriptsize] at (-1.8,1.43) {\formula};
      \node[align=center, color=uzaInk, font=\sffamily\scriptsize] at (0,-1.12) {\roots};
    \end{scope}
  }
  \begin{scope}
    \pgfmathsetmacro{\rootpos}{sqrt(0.55/0.52)}
    \draw[domain=-1.55:1.55,smooth,variable=\x,color=uzaBlue,thick] plot ({\x},{0.52*\x*\x-0.55});
    \fill[color=uzaWarn] (-\rootpos,0) circle (2pt);
    \fill[color=uzaWarn] (\rootpos,0) circle (2pt);
  \end{scope}
  \begin{scope}[xshift=5cm]
    \draw[domain=-1.55:1.55,smooth,variable=\x,color=uzaBlue,thick] plot ({\x},{0.52*\x*\x});
    \fill[color=uzaWarn] (0,0) circle (2pt);
  \end{scope}
  \begin{scope}[xshift=10cm]
    \draw[domain=-1.55:1.55,smooth,variable=\x,color=uzaBlue,thick] plot ({\x},{0.42*\x*\x+0.45});
  \end{scope}
\end{tikzpicture}
\caption{Relación entre el discriminante y el número de raíces reales de una ecuación de segundo grado.}\label{fig-discriminante}
\end{figure}

\section{Resolución de la práctica}\label{resolucion-de-la-practica}

Para resolver esta práctica se propone implementar \textbf{tres clases en Java}, cada una correspondiente a un nivel de resolución de la ecuación cuadrática.

La solución única o doble de la ecuación se calcula mediante la siguiente fórmula, conocida como la \textbf{fórmula general}:

\[
x = \frac{-b \pm \sqrt{b^2 - 4ac}}{2a} = \frac{-b \pm \sqrt{\Delta}}{2a}
\]

El valor dentro de la raíz cuadrada, \(b^2 - 4ac\), se llama \textbf{discriminante} (\(\Delta\)) y determina el número de soluciones reales que tiene la ecuación.

\begin{PracticeTipBox}{Clases en Java}
Recuerda que una clase en Java toma el siguiente formato general:

\PracticeCode{../resources/practica3/snippets/java/class-template.java}
\end{PracticeTipBox}

\subsection{Clase 1: dos soluciones reales}\label{clase-1-dos-soluciones-reales}

Si \(\Delta > 0\), la ecuación tiene dos soluciones reales distintas. Estas soluciones se calculan con la fórmula general utilizando el signo \(+\) para una raíz y el signo \(-\) para la otra.

\begin{PracticeNoteBox}{Notas para resolver el problema}
\begin{itemize}
\tightlist
\item Supondremos que solo hay raíces reales.
\item Lee los coeficientes del polinomio.
\item Aplica las fórmulas para calcular las dos soluciones.
\item Escribe las soluciones por pantalla.
\end{itemize}
\end{PracticeNoteBox}

\begin{PracticeNoteBox}{Leer datos en Java}
Utiliza la clase \texttt{Scanner} para leer los coeficientes del polinomio desde la entrada estándar:

\PracticeCode{../resources/practica3/snippets/java/scanner-read-coefficients.java}
\end{PracticeNoteBox}

\begin{PracticeNoteBox}{Imprimir resultados en Java}
Utiliza la consola para mostrar los resultados:

\PracticeCode{../resources/practica3/snippets/java/print-solution.java}
\end{PracticeNoteBox}

A continuación se muestra una posible solución en pseudocódigo:

\PracticeCode{../resources/practica3/snippets/pseudocode/raices1.txt}

\begin{PracticeTipBox}{Comprobación}
Verifica la solución con los siguientes coeficientes:

\[
a=1,\qquad b=-3,\qquad c=2
\]

El resultado esperado son dos soluciones reales: \(x=1\) y \(x=2\).
\end{PracticeTipBox}

\subsection{Clase 2: una solución real o dos soluciones reales}\label{clase-2-una-solucion-real-o-dos-soluciones-reales}

Si \(\Delta = 0\), la ecuación tiene una única solución real. En este caso, los signos \(+\) y \(-\) de la fórmula general producen el mismo resultado.

\begin{PracticeNoteBox}{Notas para resolver el problema}
\begin{itemize}
\tightlist
\item Calcula las raíces cuando \(\Delta \geq 0\).
\item Si \(\Delta > 0\), escribe las dos soluciones reales.
\item Si \(\Delta = 0\), escribe la solución única.
\item Si \(\Delta < 0\), escribe un mensaje indicando que no hay raíces reales.
\end{itemize}
\end{PracticeNoteBox}

\PracticeCode{../resources/practica3/snippets/pseudocode/raices2.txt}

\begin{PracticeTipBox}{Comprobación}
Verifica la solución con los siguientes coeficientes:

\[
a=1,\qquad b=-2,\qquad c=1
\]

El resultado esperado es una única solución real: \(x=1\).
\end{PracticeTipBox}

\subsection{Clase 3: solución completa}\label{clase-3-solucion-completa}

Si \(\Delta < 0\), la ecuación no tiene soluciones reales, ya que la raíz cuadrada de un número negativo no es un número real.

En este caso, se pueden calcular las soluciones complejas utilizando números imaginarios:

\[
x = \frac{-b \pm i\sqrt{-(b^2 - 4ac)}}{2a} = \frac{-b \pm i\sqrt{4ac - b^2}}{2a}
\]

Las dos soluciones complejas conjugadas se pueden expresar como:

\[
\begin{aligned}
x_1 &= \frac{-b + i\sqrt{4ac - b^2}}{2a} = \frac{-b}{2a} + i\frac{\sqrt{4ac - b^2}}{2a} \\
x_2 &= \frac{-b - i\sqrt{4ac - b^2}}{2a} = \frac{-b}{2a} - i\frac{\sqrt{4ac - b^2}}{2a}
\end{aligned}
\]

\begin{PracticeNoteBox}{Notas para resolver el problema}
\begin{itemize}
\tightlist
\item Calcula las raíces en todos los casos.
\item Si \(\Delta > 0\), escribe las dos soluciones reales.
\item Si \(\Delta = 0\), escribe la solución única.
\item Si \(\Delta < 0\), escribe las dos soluciones complejas conjugadas.
\end{itemize}
\end{PracticeNoteBox}

\PracticeCode{../resources/practica3/snippets/pseudocode/raices3.txt}

\begin{PracticeTipBox}{Comprobación}
Verifica la solución con los siguientes coeficientes:

\[
a=1,\qquad b=0,\qquad c=1
\]

Las soluciones esperadas son:

\[
x_1 = i,\qquad x_2 = -i
\]
\end{PracticeTipBox}

\section{Entrega de la práctica}\label{entrega-de-la-practica}

La práctica se entregará a través de la plataforma de la asignatura, con un fichero \textbf{.zip} que contenga las tres clases Java implementadas.

\begin{itemize}
\tightlist
\item \texttt{Solucion1.java} para el caso de dos soluciones reales.
\item \texttt{Solucion2.java} para el caso de discriminante no negativo, es decir, una o dos soluciones reales.
\item \texttt{Solucion3.java} para el caso completo, con soluciones reales o complejas.
\end{itemize}

\textbf{Solo realizará la entrega de la práctica un miembro del grupo.}

\end{document}
